\section{Conclusion}

This study presented a build-aligned static-runtime evidence model for evaluating the privacy and security risk posture of Android consumer apps. To address rapid app-update churn, the methodology applied a fixed 14-day evidence window as a provenance and build-alignment control, and preserved package versions, APK hashes, static and dynamic run identifiers, PCAP records, collection timestamps, and quality-assurance results. At the cutoff, all 15 apps had usable build-backed evidence, with no unresolved evidence gaps. Static analysis identified 1,235 high- and medium-severity findings and 955 unguarded-export findings across the selected builds.

Runtime analysis examined 123 captures, including 44 strict-idle, 23 quiescent foreground, and 56 interactive runs. The median additive-smoothed $\log_{2}$ interactive-to-baseline packet-rate ratio was 3.39. The monotonic association between static finding volume and interactive retained-host count was near zero (Table~\ref{tab:statistical-summary}). Together, the evidence layers provide an auditable characterization of app-level privacy and security risk posture while preserving the distinction between potential exposure and observed runtime behavior; neither layer alone should be read as exploitability, malicious activity, or a calibrated probability of harm.

Future work will extend this evidence model through longitudinal build-to-build comparisons, additional devices and Android versions, improved service-family labeling, and stronger held-out runtime validation. Applying the same build-aligned evidence contract across successive collection windows would support comparisons of app updates while preserving provenance and the study's privacy boundaries.
